%!TEX program = xelatex
% =============================================================================
% GUION DE LA DEFENSA — Caso 2: Línea de Producción con Balanceo Dinámico
%
% Versión imprimible del guion hablado, para sostener en la mano durante la
% defensa.
%
% OJO: el mismo texto existe además en docs/guion-defensa.md, para leerlo en
% GitHub. Son dos copias del mismo contenido y pueden desincronizarse. Este
% archivo es el que manda, porque es el que produce el PDF: si hay que cambiar
% algo, se cambia acá, se recompila, y después se refleja en el .md.
%
% Compilar: xelatex guion-defensa.tex   (una pasada basta)
% =============================================================================
\documentclass[11pt,a4paper]{article}

\usepackage[spanish]{babel}
\deactivatequoting
\usepackage[margin=2.4cm, top=2.2cm, bottom=2.4cm]{geometry}
\usepackage{fontspec}
\usepackage{xcolor}
\usepackage{enumitem}
\usepackage{titlesec}
\usepackage{fancyhdr}
\usepackage{parskip}
\usepackage{setspace}

% Misma cascada de fuentes que el resto del proyecto.
\IfFontExistsTF{Segoe UI}
  {\setsansfont{Segoe UI}}
  {\IfFontExistsTF{Lato-Regular.ttf}
    {\setsansfont{Lato}[Extension=.ttf, UprightFont=*-Regular, BoldFont=*-Bold,
       ItalicFont=*-Italic, BoldItalicFont=*-BoldItalic]}{}}
\renewcommand{\familydefault}{\sfdefault}

\definecolor{tinta}{HTML}{1C2026}
\definecolor{tinta2}{HTML}{52514E}
\definecolor{suave}{HTML}{898781}
\definecolor{tealu}{HTML}{009B8A}
\definecolor{ambar}{HTML}{B37A0F}

\color{tinta}
\setstretch{1.22}

% Que ninguna palabra quede partida entre dos páginas. En un guion que se lee en
% voz alta, tener que dar vuelta la hoja a mitad de palabra es un tropiezo.
\brokenpenalty=10000

\pagestyle{fancy}
\fancyhf{}
\fancyhead[L]{\footnotesize\color{suave}Guion de la defensa · Caso 2 — Línea de producción con balanceo dinámico}
\fancyhead[R]{\footnotesize\color{suave}\thepage}
\renewcommand{\headrulewidth}{0pt}

\titleformat{\section}{\large\bfseries\color{tealu}}{}{0pt}{}
\titlespacing*{\section}{0pt}{16pt}{6pt}

% Encabezado de lámina: (Diapositiva N — Título)
\newcommand{\lamina}[1]{%
  \vspace{12pt}%
  \noindent{\bfseries\color{tealu}(#1)}\par\vspace{3pt}}

% Texto hablado: filete de acento a la izquierda.
\newenvironment{decir}
  {\begin{list}{}{\setlength{\leftmargin}{14pt}\setlength{\rightmargin}{0pt}%
    \setlength{\itemindent}{0pt}\setlength{\listparindent}{0pt}%
    \setlength{\parsep}{6pt}}\item[]%
    \color{tinta}}
  {\end{list}}

\newcommand{\opcional}{{\bfseries\color{ambar}$\langle$opcional$\rangle$}\ }

\begin{document}

% ── PORTADA ──────────────────────────────────────────────────────────────────
\thispagestyle{empty}
\vspace*{1.2cm}
{\color{tealu}\rule{3cm}{2pt}}\\[14pt]
{\fontsize{24}{29}\selectfont\bfseries Guion de la defensa}\\[8pt]
{\large\color{tinta2}Caso 2 — Línea de producción con balanceo dinámico}\\[6pt]
{\color{tinta2}Estructura de Computadores y Sistemas Distribuidos · Evaluación 3}\\[4pt]
{\color{suave}César Olivares · Simón García-Huidobro · 26 de agosto de 2026}

\vspace{18pt}
{\color{suave}\rule{\linewidth}{0.4pt}}
\vspace{10pt}

\noindent Guion hablado para las 21 láminas de
\texttt{presentacion\_caso2\_alt.pdf}. Está escrito para decirse en voz alta,
no para leerse en pantalla.

\vspace{6pt}
\noindent\textbf{Duración.} El guion tiene \textbf{1.762 palabras} habladas. A
145 palabras por minuto, que es un ritmo de exposición normal, son
\textbf{12 minutos}. Con 4 de demo da \textbf{16}, o sea \textbf{uno de más}.
Hay que ganar ese minuto de una de estas dos formas, y decidirlo en el ensayo,
no en la defensa:

\begin{itemize}[leftmargin=16pt, itemsep=2pt, topsep=4pt]
  \item \textbf{Bajar la demo a 3 minutos.} Es lo más fácil: el paso 5 (matar
        Redis) se puede contar en vez de hacerse, porque su evidencia está en
        los anexos del informe.
  \item \textbf{Cortar el cuarto acto} (láminas 18 y 19, Resiliencia). Es el
        contenido menos exigido por el caso y sale sin romper el hilo: de la
        lámina 17 se pasa directo a conclusiones.
\end{itemize}

\noindent Los párrafos marcados con \opcional son los primeros que deben caer
si el cronómetro aprieta. La regla es cortar contenido, nunca hablar más
rápido.

\vspace{6pt}
\noindent\textbf{Turnos.} Los cuatro actos son los cortes naturales. Una
repartición pareja es Acto I + Acto III para uno, Acto II + Acto IV para el
otro, y la demo entre los dos. Definirlo antes y anotarlo acá:
\underline{\hspace{5cm}}

\newpage

% ── GUION ────────────────────────────────────────────────────────────────────
\lamina{Diapositiva 1 — Portada}
\begin{decir}
``Hola, profesor. Hoy presentamos el Caso 2: una línea de producción con
balanceo dinámico de carga. Construimos un sistema distribuido que simula una
planta conservera, detecta en vivo dónde se está atascando la producción, y
responde una pregunta concreta: qué pasa con ese atasco cuando uno agrega más
máquinas a la etapa lenta.''
\end{decir}

\lamina{Diapositivas 2 y 3 — Primer acto: El problema}
\begin{decir}
``La planta tiene cuatro etapas en serie: fileteado, envasado, sellado y
esterilización. La unidad de trabajo es el \textbf{lote}, una cantidad de latas
del mismo formato, y cada lote pasa por las cuatro en orden.

Cada etapa demora algo distinto por lote: cuatro segundos, doce, tres y siete.
A ese número lo llamamos \textbf{tiempo de ciclo}, y es una propiedad de la
máquina: la autoclave demora siete segundos y no hay cómo apurarla.

Entre etapa y etapa hay una \textbf{cola}, que es la fila de lotes esperando su
turno. Y acá está el concepto central: si a una etapa le llega trabajo más
rápido de lo que alcanza a procesarlo, su fila crece sin parar. Eso es un
\textbf{cuello de botella}.

Envasado es la etapa más lenta, así que es la que se \textbf{replica}: levantar
varias copias idénticas funcionando en paralelo para que se repartan el
trabajo.

La pregunta del caso es entonces: ¿qué le pasa al cuello de botella cuando
escalamos la etapa lenta? Les adelanto la respuesta, porque todo lo que viene
es la evidencia de esta frase: \textbf{el cuello no desaparece, se muda}.''
\end{decir}

\lamina{Diapositivas 4 y 5 — Segundo acto: Arquitectura}
\begin{decir}
``Usamos cinco \textbf{servicios}. Un servicio es un programa que corre por su
cuenta y le ofrece algo a los demás por la red, en vez de ser una función
dentro del mismo programa.

Los agrupamos en tres \textbf{nodos lógicos}, y la agrupación no es arbitraria:
cada nodo junta lo que escala o falla por el mismo motivo. El A le da la cara
al operador; el B tiene las cuatro estaciones, que es lo único que se
multiplica bajo carga; y el C concentra el estado compartido, que por ser
compartido no se puede duplicar sin coordinación. Separar nodos separa dominios
de falla: si las estaciones saturan el procesador, el operador sigue siendo
atendido.

Un detalle del que estamos conformes: las cuatro estaciones son \textbf{el
mismo programa}. No hay un código de fileteado y otro de envasado; la etapa que
le toca y su tiempo de ciclo le llegan por configuración al arrancar. Todo
corre sobre \textbf{Docker}, que empaqueta cada programa con lo que necesita en
una unidad llamada contenedor, de modo que funciona igual en cualquier máquina.
Por eso levantar el sistema completo es un comando y no una lista de pasos de
instalación.''
\end{decir}

\lamina{Diapositiva 6 — Balanceo por demanda}
\begin{decir}
``¿Cómo se reparte el trabajo entre las réplicas?

Lo habitual es poner un \textbf{balanceador} delante: una pieza que recibe cada
lote y decide a cuál réplica mandárselo. La forma más simple es
\textbf{round-robin}: por turnos fijos y en orden, uno para cada una.

Nosotros lo hicimos al revés, y es una decisión que defendemos. En vez de que
alguien les asigne trabajo, las réplicas \textbf{piden}: cada una toma el
siguiente lote de la cola compartida cuando queda libre. Nuestro balanceador es
la cola.

¿Por qué es mejor? Porque round-robin es estático: le sigue mandando lotes a
una réplica lenta al mismo ritmo que a las rápidas. Pidiendo, la réplica lenta
simplemente pide menos veces. El balanceo se adapta solo, y más adelante
mostramos el experimento que lo mide.''
\end{decir}

\lamina{Diapositiva 7 — Condición de carrera}
\begin{decir}
``Pero pedir de una cola compartida abre un riesgo, y es el problema central
del caso.

Si dos réplicas piden al mismo tiempo, el lote se lo tiene que llevar
exactamente una. Cuando el resultado depende de cuál llegó primero, eso es una
\textbf{condición de carrera}.

La provocamos a propósito antes de resolverla, y las dos versiones quedaron en
el historial de commits. La versión ingenua reclama en dos pasos: primero mira
cuál es el próximo lote, después lo saca. Entre el paso uno y el paso dos el
lote sigue en la cola, así que otra réplica que mire en ese instante ve el
mismo. Ensanchamos esa ventana hasta reproducir duplicados: con tres réplicas y
cien pedidos, la planta envasaba trescientos lotes. El triple de materia prima.

La solución no fue ponerle un candado con expiración, que es lo clásico. Usamos
una \textbf{operación atómica}: una operación que ocurre entera o no ocurre,
sin ningún instante intermedio en que otro pueda ver el sistema a medio hacer.
\textbf{Redis}, la base de datos en memoria que usamos como cola, atiende sus
comandos de uno en uno, así que sacar un lote es indivisible por construcción.

La diferencia conceptual es la que nos importa: un candado \textbf{administra}
la ventana de riesgo; la operación atómica la \textbf{elimina}. Con ella
medimos cero duplicados en doscientas órdenes.''
\end{decir}

\lamina{Diapositiva 8 — El criterio de cuello de botella}
\begin{decir}
``Antes de los resultados hay que definir cómo decidimos dónde está el atasco,
porque no hay una definición única y es la decisión técnica más importante que
tomamos.

Hay dos tiempos que distinguir. El \textbf{tiempo de servicio} es lo que la
etapa demora procesando un lote que ya tiene en la mano. El \textbf{tiempo de
espera} es lo que el lote pasa haciendo fila antes de que lo tomen.

La analogía es el supermercado: un cajero lento no es un atasco. El atasco es
la fila que crece frente a su caja.

Nuestro criterio mide \textbf{espera, no servicio}. El cuello es la etapa con
mayor espera promedio sobre una \textbf{ventana móvil} de sesenta segundos
---mirando solo el último minuto, para que una congestión vieja y ya resuelta
no contamine el diagnóstico de ahora--- comparada contra umbrales relativos al
ciclo de cada etapa: vez y media es advertencia, tres veces es crítico.

Y no es una sutileza. Con dos réplicas, envasado sigue teniendo el ciclo más
largo de la línea y \textbf{no} es el cuello. Midiendo servicio, el detector
habría señalado envasado para siempre y nunca habríamos visto el fenómeno que
veníamos a estudiar.''
\end{decir}

\lamina{Diapositivas 9 y 10 — Demo en vivo, $\sim$4 minutos}
\begin{decir}
``Vamos a mostrarlo funcionando. Esto está corriendo de verdad: dos réplicas de
envasado y las otras tres etapas.''
\end{decir}

\noindent\textbf{Guion de la demo:}
\begin{enumerate}[leftmargin=18pt, itemsep=4pt, topsep=4pt]
  \item \textbf{La línea al día.} Mostrar las cuatro etapas, las réplicas con su
        ciclo y su carga, y el cartel verde.
  \item \textbf{Ingresar lotes} desde el panel izquierdo, \textbf{espaciados
        unos cinco segundos}.

        {\small\color{ambar}\textbf{Ojo:} no inyectar muchos de golpe. Se apilan
        en la primera cola y el tablero marca fileteado como el atasco ---
        correctamente, pero es la respuesta a otra pregunta. Si alguien lo
        nota, explicar exactamente eso.}
  \item \textbf{Ver crecer la espera} de esterilización. El cartel pasa de verde
        a ámbar y luego a rojo.

        {\small\itshape ``Fíjense que el cartel no dice solo `hay un problema':
        nombra la etapa, dice cuántos lotes hacen cola y da el número. Ese
        diagnóstico viene del servicio de métricas; el tablero solo lo
        muestra.''}
  \item \textbf{La condición de carrera, en vivo.} Bajar y apretar «Ejecutar la
        comparación».

        {\small\itshape ``Esto corre el experimento ahora: encola los mismos
        lotes dos veces, una con cada versión del reclamo. Con cien pedidos, la
        ingenua envasa trescientos y la atómica cien. Y es \textbf{el mismo
        código} que corre en las estaciones: las dos funciones se importan del
        mismo módulo, no hay una copia para la demo que pueda quedar
        desactualizada. Usa una cola aparte, así que la línea de arriba sigue
        produciendo.''}
  \item \textbf{Matar Redis.} Detener el contenedor, mostrar que la interfaz
        nombra al servicio caído sin un error críptico, y volver a levantarlo.
\end{enumerate}

\noindent{\small\color{suave}\itshape Plan B si algo falla: las capturas están
en los anexos del informe.}

\lamina{Diapositivas 11 y 12 — Tercer acto: la aritmética del escenario}
\begin{decir}
``Los experimentos. Todos reproducibles y con sus datos crudos versionados.

El escenario es siempre el mismo: una orden cada cinco segundos durante tres
minutos, o sea doce lotes por minuto entrando.

Este gráfico es la aritmética de capacidad, y lo mostramos \textbf{antes} de
los resultados a propósito, porque los predice. Cada barra es cuántos lotes por
minuto puede sacar cada etapa; la línea punteada son los doce que llegan. Toda
etapa bajo esa línea acumula fila tarde o temprano.

Envasado con una réplica saca cinco. Con dos sube a diez: mejor, pero todavía
insuficiente. Y ahí aparece esterilización con ocho coma seis, que queda como
el siguiente límite. La figura ya anticipa el resultado.''
\end{decir}

\lamina{Diapositiva 13 — El resultado central}
\begin{decir}
``Esto es lo que medimos: la espera promedio de las cuatro etapas, con una, dos
y tres réplicas de envasado.

Con \textbf{una} réplica, envasado acumula sesenta y siete coma ocho segundos
de espera y el resto de la línea está casi ociosa.

Con \textbf{dos}, el cuello \textbf{se desplaza a esterilización}. Y ojo con
esto: esterilización no cambió nada, sigue demorando siete segundos por lote.
Lo que cambió es que ahora le llega trabajo más rápido de lo que puede sacarlo,
porque envasado dejó de ser el freno.

Con \textbf{tres}, envasado queda sobrado, con menos de un segundo de espera\dots{}
y el cuello sigue en esterilización. Ahí está la respuesta: el cuello no se
eliminó. Se mudó.''
\end{decir}

\lamina{Diapositiva 14 — Cuánto compra cada réplica}
\begin{decir}
``La pregunta del jefe de planta es cuánto le compró esa inversión. Para eso
está el \textbf{lead time}: el tiempo total de punta a punta, desde que se crea
la orden hasta que el lote sale terminado. Suma todas las esperas y todos los
ciclos del recorrido.

La segunda réplica compra un veintiuno por ciento de mejora. La tercera,
prácticamente nada.

Y esa es la conclusión útil: la siguiente inversión correcta no es una tercera
envasadora, es \textbf{una segunda autoclave}. El tablero responde eso en vivo,
que es para lo que sirve medir.''
\end{decir}

\lamina{Diapositiva 15 — Balanceo con una réplica lenta}
\begin{decir}
``Este experimento prueba que el balanceo es dinámico de verdad.

Tres réplicas de envasado sobre la misma cola, pero una configurada al doble de
lenta, y treinta lotes.

Un round-robin habría entregado diez, diez y diez sin mirar la velocidad de
nadie. Lo que medimos fue \textbf{once, doce y siete}: la lenta procesó un
cuarenta por ciento menos.

Nadie programó esa proporción, profesor. No hay una línea de código que diga
`dale menos a la lenta'. Salió sola de que cada réplica pide trabajo únicamente
cuando queda libre. Eso es lo que hace que el balanceo sea dinámico:
proporcional a la capacidad real de cada máquina, no a un turno
preestablecido.''
\end{decir}

\lamina{Diapositiva 16 — Saturación y recuperación}
\begin{decir}
``El último experimento fue meterle una ráfaga de veinticinco órdenes de golpe.

A los treinta y dos segundos el sistema \textbf{se declara crítico solo}.

\opcional Y se ve algo bonito: la congestión recorre la línea como una ola,
aparece en la etapa de entrada y se propaga aguas abajo.

A los doscientos setenta y ocho vuelve todo a normal \textbf{sin intervención}.
Es consecuencia directa de la ventana móvil: las esperas viejas salieron del
último minuto. El sistema no se `arregló'; dejó de estar congestionado y la
métrica lo reflejó.''
\end{decir}

\lamina{Diapositiva 17 — Robustez}
\begin{decir}
``Tres cosas que sostienen todo lo demás.

\textbf{Errores.} Al matar Redis, el servicio de órdenes responde en dos
segundos nombrando al servicio que falta. \opcional Antes se colgaba cuarenta y
seis: le habíamos puesto un límite de tiempo, pero el cuelgue estaba en la
resolución del nombre y no en la conexión, y ese límite no lo cubría.

\textbf{Persistencia.} Guardamos en \textbf{SQLite}, una base de datos que vive
en un solo archivo. Dos tablas: órdenes y eventos. Los tiempos \textbf{no se
guardan}, se calculan desde los eventos --- el dato guardado dos veces es el
que termina contradiciéndose.

\textbf{Despliegue.} Un comando levanta todo desde cero, probado clonando en
una carpeta vacía: ocho de ocho contenedores sanos.''
\end{decir}

\lamina{Diapositivas 18 y 19 — Cuarto acto: Resiliencia}
\begin{decir}
``El caso no pedía tolerancia a fallos y deliberadamente no la implementamos:
el sistema \textbf{detecta y explica} las fallas en vez de disimularlas. Pero
sabemos dónde iría cada pieza.

Las \textbf{estaciones} no tienen estado propio: mueren sin corromper nada.
Falta re-encolar el lote en curso si una muere a la mitad. \textbf{Redis} es el
punto único de falla del reparto: ahí iría una réplica de respaldo. Y
\textbf{SQLite} aguanta esta carga, pero con más escritores habría que migrar, y
eso toca un solo módulo.

Lo que queremos destacar: ninguna de esas mejoras obliga a rediseñar las otras.
Para eso servía haber dividido el sistema en nodos desde el principio.''
\end{decir}

\lamina{Diapositivas 20 y 21 — Conclusiones y cierre}
\begin{decir}
``Para cerrar, cinco frases, todas medidas.

\textbf{El cuello de botella no se elimina, se muda:} de envasado a
esterilización al pasar de una a dos réplicas.

\textbf{Medir espera en vez de servicio es lo que permite verlo:} el cuello
real tenía un ciclo más corto que envasado.

\textbf{El balanceo dinámico emerge del mecanismo:} once, doce y siete con una
réplica lenta, sin que nadie asignara ese reparto.

\textbf{La condición de carrera se elimina, no se administra:} cero duplicados
en doscientas órdenes.

\textbf{Las fallas se explican, no se enmascaran:} un error claro en dos
segundos, en vez de cuarenta y seis de silencio.

Si me quedo con una sola idea de todo el proyecto, profesor, es esta:
\textbf{agregar máquinas no elimina el cuello de botella, lo traslada}. Por eso
lo que hay que construir no es una línea más rápida, sino una línea que sepa
decirte dónde se está atascando ahora mismo. Sin eso, uno invierte a ciegas.

Muchas gracias.''
\end{decir}

% ── PREGUNTAS PROBABLES ──────────────────────────────────────────────────────
% Sin salto de página forzado: dejarlo abría una página casi en blanco. El
% título en color ya separa lo suficiente.
\vspace{10pt}
\section{Preguntas probables}

\noindent\textbf{¿Dónde está el balanceador?}\\
Es la cola compartida: el punto único por el que pasa todo el trabajo y que
decide quién procesa qué. Elegimos reparto por demanda porque un balanceador
clásico no es dinámico --- le sigue mandando trabajo a la réplica lenta. El
experimento de la réplica lenta lo muestra adaptándose solo.

\vspace{8pt}
\noindent\textbf{¿Y el lock o contador distribuido?}\\
Lo cubre el mismo mecanismo. La extracción atómica de Redis es un candado
implícito por elemento, porque Redis atiende sus comandos de uno en uno. Un
candado explícito con expiración habría \emph{administrado} la ventana de
carrera; la extracción atómica la \emph{elimina}. Y lo demostramos: dos pasos
da duplicados, un paso da cero en doscientas.

\vspace{8pt}
\noindent\textbf{¿Por qué el cuello es esterilización si envasado es más lento
(12 s contra 7 s)?}\\
Porque con dos réplicas el tiempo efectivo de envasado es de unos seis segundos
por lote. Y en general el cuello no es la etapa de mayor ciclo sino la de mayor
espera: la que recibe trabajo más rápido de lo que puede drenarlo.

\vspace{8pt}
\noindent\textbf{¿Qué pasa si una réplica muere a mitad de un lote?}\\
Ese lote ya salió de la cola y se pierde de la línea. Es una limitación conocida
y aceptada, porque el caso excluía tolerancia a fallos. Sabemos el camino: mover
el lote a una cola «en proceso» en vez de sacarlo, y un barrendero que lo
devuelva al detectar que la réplica dejó de dar señales.

\vspace{8pt}
\noindent\textbf{¿Por qué SQLite y no un motor con servidor?}\\
Un solo nodo de estado, escrituras cortas, y cero infraestructura extra que
administrar en la demo. Lo estresamos en el experimento de la réplica lenta y
encontramos su límite, que está documentado. El cambio a otro motor toca un solo
módulo si hiciera falta.

\vspace{8pt}
\noindent\textbf{¿Cómo saben que las métricas son correctas?}\\
Los tiempos de servicio que medimos coinciden con los ciclos configurados
---cuatro coma ocho, doce coma cuatro, tres coma cuatro y siete coma tres,
contra cuatro, doce, tres y siete---. La instrumentación se valida sola. Y todo
se calcula desde la tabla de eventos, así que no hay un segundo lugar que pueda
contradecirla.

\vspace{8pt}
\noindent\textbf{¿Por qué la espera se mide contra la entrada a cada cola y no
contra la creación de la orden?}\\
Porque con cuatro etapas en serie, medir contra la creación le cargaría a
sellado todo lo que pasó antes en fileteado y envasado, y señalaría al culpable
equivocado. Cada etapa responde solo por su propia cola.

\end{document}
